\documentclass[12pt,a4paper]{article}
\usepackage[utf8]{vietnam}
\usepackage{amsmath,amssymb}
\usepackage{geometry}
\geometry{
  top=12mm,
  bottom=12mm,
  left=14mm,
  right=14mm,
  footskip=6mm,
  headheight=14pt
}

\usepackage{tikz}
\usetikzlibrary{arrows.meta,calc,shapes.geometric}
\usepackage[most]{tcolorbox}
\usepackage{fontawesome}
\usepackage{fancyhdr}
\usepackage{enumitem}

\pagestyle{fancy}
\fancyhf{}
\fancyhead[L]{\footnotesize\color{teal!80!black}\faBook\ \textbf{SỔ TAY KIẾN THỨC TOÁN 3}}
\fancyhead[R]{\footnotesize\color{teal!80!black}\textbf{Lớp Toán Thầy Lĩnh -- loptoanthaylinh.vn}}
\fancyfoot[C]{\footnotesize\color{gray}-- Sổ tay kiến thức: Bài 20. Thực hành vẽ hình và trang trí --}
\renewcommand{\headrulewidth}{0.4pt}
\renewcommand{\footrulewidth}{0pt}

\setlength{\parindent}{0pt}
\setlength{\parskip}{0pt}

\begin{document}

% --- BANNER BÀI HỌC ---
\begin{tcolorbox}[
  enhanced,
  colback=teal!8!white,
  colframe=teal!70!black,
  arc=3mm, boxrule=1.2pt,
  left=3mm, right=3mm, top=2mm, bottom=2mm,
  halign=center
]
  {\LARGE\bfseries\color{teal!90!black}\faBookmark\ BÀI 20: THỰC HÀNH VẼ HÌNH VÀ TRANG TRÍ}\\[1mm]
  {\itshape\footnotesize\color{black!75} Kĩ thuật vẽ trên giấy kẻ ô vuông, ghép nối hình học sáng tạo và trang trí đối xứng trên 1 trang A4}
\end{tcolorbox}

\vspace{-2mm}

\noindent
\begin{minipage}[t]{0.49\textwidth}\vspace{0pt}
\begin{tcolorbox}[
  enhanced,
  colback=blue!4!white,
  colframe=blue!75!black,
  arc=2mm, boxrule=0.8pt,
  left=2mm, right=2mm, top=1.8mm, bottom=1.8mm,
  title={\footnotesize\bfseries\color{white}\faPencil\ 1. Kĩ thuật vẽ trên lưới ô vuông}
]
{\footnotesize
\begin{itemize}\setlength{\itemsep}{1pt}\setlength{\leftskip}{-2.5mm}
  \item \textbf{Lợi ích của lưới ô vuông:} Giúp căn chỉnh góc vuông và đếm độ dài đoạn thẳng cực kì chính xác.
  \item \textbf{Các bước vẽ hình chữ nhật, hình vuông:}
  \begin{enumerate}\setlength{\itemsep}{0.5pt}\setlength{\leftskip}{-2.5mm}
    \item Xác định $4$ đỉnh theo số ô vuông quy định.
    \item Dùng thước thẳng nối từng cặp đỉnh lại với nhau.
    \item Kiểm tra lại số ô dài, rộng và $4$ góc vuông.
  \end{enumerate}
\end{itemize}
}
\end{tcolorbox}
\end{minipage}
\hfill
\begin{minipage}[t]{0.49\textwidth}\vspace{0pt}
\begin{tcolorbox}[
  enhanced,
  colback=teal!4!white,
  colframe=teal!75!black,
  arc=2mm, boxrule=0.8pt,
  left=2mm, right=2mm, top=1.8mm, bottom=1.8mm,
  title={\footnotesize\bfseries\color{white}\faPaintBrush\ 2. Ghép hình \& Trang trí đối xứng}
]
{\footnotesize
\begin{itemize}\setlength{\itemsep}{1pt}\setlength{\leftskip}{-2.5mm}
  \item \textbf{Ghép hình sáng tạo:}
  \begin{itemize}\setlength{\itemsep}{0.5pt}\setlength{\leftskip}{-2.5mm}
    \item Ghép $2$ tam giác vuông cân $\to$ $1$ hình vuông.
    \item Ghép $2$ hình vuông $\to$ $1$ hình chữ nhật.
  \end{itemize}
  \item \textbf{Hoa văn trang trí đối xứng:} Lặp lại các họa tiết hình học (hình tròn, vuông, thoi) tạo thành đường diềm viền quanh hoặc tâm đối xứng.
\end{itemize}
}
\end{tcolorbox}
\end{minipage}

\vspace{-2.5mm}
% --- KHỐI 3: SƠ ĐỒ TƯ DUY TIKZ ---
\noindent
\begin{tcolorbox}[
  enhanced,
  colback=orange!4!white,
  colframe=orange!80!black,
  arc=2mm, boxrule=0.8pt,
  left=2mm, right=2mm, top=1.5mm, bottom=1.5mm,
  title={\footnotesize\bfseries\color{white}\faSitemap\ SƠ ĐỒ TƯ DUY: NGHỆ THUẬT GHÉP TÁCH HÌNH HỌC TRÊN LƯỚI Ô}
]
\centering
\begin{tikzpicture}[scale=0.85]
  % Luoi o vuong minh hoa
  \draw[step=0.6cm, gray!40, very thin] (0,0) grid (3.6, 1.8);
  \draw[thick, blue!80!black, fill=blue!15] (0.6, 0.6) rectangle (3.0, 1.2);
  \node[font=\scriptsize\bfseries, blue!80!black] at (1.8, 0.9) {Dài 4 ô, Rộng 1 ô};

  \draw[thick, ->, purple!80!black] (4.0, 0.9) -- (5.0, 0.9) node[midway, above, font=\tiny\bfseries] {Ghép 2};

  \draw[thick, teal!80!black, fill=teal!20] (5.5, 0) rectangle (7.5, 1.8);
  \draw[dashed, red!80!black] (5.5, 0) -- (7.5, 1.8);
  \node[font=\scriptsize\bfseries, teal!90!black] at (6.5, 0.9) {2 Tam giác};

  \draw[thick, ->, purple!80!black] (8.0, 0.9) -- (9.0, 0.9) node[midway, above, font=\tiny\bfseries] {Xoay};

  \draw[thick, orange!80!black, fill=orange!15] (9.6, 0.9) circle (0.9cm);
  \fill[black] (9.6, 0.9) circle (1.5pt);
  \node[font=\scriptsize\bfseries, orange!90!black] at (9.6, 0.9) [below] {Tâm compa};
\end{tikzpicture}
\end{tcolorbox}

\vspace{-2.5mm}
% --- KHỐI 4: VÍ DỤ & MẸO TRÁNH BẪY ---
\noindent
\begin{minipage}[t]{0.49\textwidth}\vspace{0pt}
\begin{tcolorbox}[
  enhanced,
  colback=green!4!white,
  colframe=green!60!black,
  arc=2mm, boxrule=0.8pt,
  left=2mm, right=2mm, top=1.8mm, bottom=1.8mm,
  title={\footnotesize\bfseries\color{white}\faCheckCircle\ Ví dụ thực hành chuẩn mực}
]
{\footnotesize
\textbf{1. Vẽ hình chữ nhật trên ô vở:}\\
Vẽ hình chữ nhật có chiều dài $5$ ô, chiều rộng $3$ ô. Ta chấm $4$ đỉnh theo góc của các ô vuông rồi dùng thước nối lại.\\
\textbf{2. Đếm số hình tam giác:}\\
Vẽ một hình chữ nhật rồi vẽ cả $2$ đường chéo cắt nhau, ta sẽ đếm được tất cả $\mathbf{8}$ hình tam giác (gồm 4 tam giác nhỏ và 4 tam giác lớn).
}
\end{tcolorbox}
\end{minipage}
\hfill
\begin{minipage}[t]{0.49\textwidth}\vspace{0pt}
\begin{tcolorbox}[
  enhanced,
  colback=red!4!white,
  colframe=red!70!black,
  arc=2mm, boxrule=0.8pt,
  left=2mm, right=2mm, top=1.8mm, bottom=1.8mm,
  title={\footnotesize\bfseries\color{white}\faExclamationTriangle\ Mẹo tránh bẫy trắc nghiệm}
]
{\footnotesize
\begin{itemize}\setlength{\itemsep}{1.5pt}\setlength{\leftskip}{-2.5mm}
  \item \textbf{Bẫy vẽ méo đường thẳng:}\\
  Khi vẽ trên giấy ô vuông, phải đặt cạnh thước thẳng trùng khít lên đường kẻ in sẵn của vở.
  \item \textbf{Bẫy đếm sót hình tam giác ghép:}\\
  Học sinh chỉ đếm các hình tam giác đơn lẻ mà thường bỏ quên các tam giác lớn do 2 hoặc nhiều tam giác nhỏ ghép lại!
\end{itemize}
}
\end{tcolorbox}
\end{minipage}


\end{document}
